% Source: https://artofproblemsolving.com/community/c3h1431988_pg_252_intro_to_geo_sidenote
% Stable identity: aops-topic-1431988-q01
Let $O$ be the centre. Joining $O$ to the twelve vertices cuts the dodecagon into $12$ congruent isosceles triangles, each with two sides $R$ and apex angle $\dfrac{360^\circ}{12}=30^\circ$. Take one of them, $OAB$, and drop the perpendicular $BH$ to $OA$.

\begin{asy}
usepackage("amsmath");
size(5cm);
defaultpen(fontsize(10pt));
pair ft(pair P,pair U,pair V){return U+dot(P-U,V-U)/dot(V-U,V-U)*(V-U);}
path ra(pair P,pair u,pair v,real s=0.2){return P+s*unit(u)--P+s*unit(u)+s*unit(v)--P+s*unit(v);}
real R=2; pair O=(0,0);
pair V(int k){return R*dir(30*k);}
path g=V(0); for(int k=1;k<12;++k) g=g--V(k); g=g--cycle;
draw(g,linewidth(0.9));
draw(O--V(0),blue); draw(O--V(1),blue);
pair H=ft(V(1),O,V(0));
draw(V(1)--H,red+dashed); draw(ra(H,O-H,V(1)-H,0.12),red);
draw(arc(O,0.45,0,30)); label("$30^\circ$",0.75*dir(15),fontsize(8pt));
dot(O);dot(V(0));dot(V(1));dot(H,red);
label("$O$",O,W); label("$A$",V(0),E); label("$B$",V(1),NE); label("$H$",H,S,red);
label("$R$",(O+V(1))/2,NW,blue);
label("$\dfrac{R}{2}$",(V(1)+H)/2,W,red);
label("\mbox{}",(min(currentpicture,user=true)+max(currentpicture,user=true))/2,black); // keeps the text after the figure black when LaTeX typesets the labels inline
\end{asy}

Triangle $OHB$ is a $30$-$60$-$90$ triangle with hypotenuse $OB=R$, so $BH=\dfrac{R}{2}$. Hence
\[
[OAB]=\dfrac{1}{2}\cdot OA\cdot BH=\dfrac{1}{2}\cdot R\cdot\dfrac{R}{2}=\dfrac{R^2}{4},
\]
and the dodecagon has area $12\cdot\dfrac{R^2}{4}=\boxed{3R^2}$. $\blacksquare$

(The dissection in the Art of Problem Solving logo rearranges the pieces of the dodecagon into three squares of side $R$, which shows the same fact.)
